\documentclass[11pt]{article}
\usepackage[margin=1in]{geometry}
\usepackage{xcolor}
\usepackage{enumitem}
\usepackage{hyperref}
\usepackage{titlesec}

% ============================================================
% Team comment commands — usage: \yiling{your comment here}
% Each renders as [Name: comment] in that member's color.
% ============================================================
\newcommand{\yiling}[1]{\textcolor{red}{[Yiling: #1]}}
\newcommand{\vincent}[1]{\textcolor{blue}{[Vincent: #1]}}
\newcommand{\haoran}[1]{\textcolor{teal}{[Haoran: #1]}}
\newcommand{\jenny}[1]{\textcolor{violet}{[Jenny: #1]}}
\newcommand{\haojun}[1]{\textcolor{orange}{[Haojun: #1]}}

\title{Project Proposal: [Project Name]}
\author{[Team Name]}
\date{CPSC 4390/5390 --- Fall 2026}

\begin{document}
\maketitle

% \yiling{example comment -- delete once everyone knows how this works}

\section{Team}

\textbf{Project / Team Name:} [TBD]

\textbf{Team Members:}
\begin{itemize}[nosep]
    \item Yiling --- [role / year]
    \item Vincent --- [role / year]
    \item Haoran --- [role / year]
    \item Jenny --- [role / year]
    \item Haojun --- [role / year]
\end{itemize}

\textbf{Point of contact with instructors:} [Name]

\section{Problem Statement}
% Describe what problem the system solves and why it matters / is exciting.
[...]

\section{Target Audience}
% Who will use this system? Be specific (e.g., not "students" but "X students who do Y").
[...]

\section{System Overview \& Expected Features}
% High-level description of what the system does, then a bullet list of features.
[...]

\subsection{Feature List}
\begin{itemize}
    \item [Feature 1]
    \item [Feature 2]
    \item [Feature 3]
\end{itemize}

\section{Minimum Viable Product (MVP)}
% Which subset of the features above is the MVP -- the smallest useful version.
[...]

\section{Scope Justification}
% Argue the project is neither too easy nor too ambitious.
% Prioritize features so lower-priority ones can be dropped if time runs short.
% If the project is ambitious, sketch the major tasks so scope is credible.
[...]

\section{Technology Stack}
% Languages, frameworks, tools -- and WHY each was chosen.
\begin{itemize}
    \item \textbf{Language(s):} [...]
    \item \textbf{Framework(s):} [...]
    \item \textbf{Tools} (CI/CD, deployment, testing, etc.): [...]
\end{itemize}

\section{Team Background}
% For each member: languages + expertise level, largest solo project, largest team project (+ team size).

\subsection{Yiling}
\begin{itemize}[nosep]
    \item Languages / expertise: [...]
    \item Largest individual project: [...]
    \item Largest team project (team size): [...]
\end{itemize}

\subsection{Vincent}
\begin{itemize}[nosep]
    \item Languages / expertise: [...]
    \item Largest individual project: [...]
    \item Largest team project (team size): [...]
\end{itemize}

\subsection{Haoran}
\begin{itemize}[nosep]
    \item Languages / expertise: [...]
    \item Largest individual project: [...]
    \item Largest team project (team size): [...]
\end{itemize}

\subsection{Jenny}
\begin{itemize}[nosep]
    \item Languages / expertise: [...]
    \item Largest individual project: [...]
    \item Largest team project (team size): [...]
\end{itemize}

\subsection{Haojun}
\begin{itemize}[nosep]
    \item Languages / expertise: [...]
    \item Largest individual project: [...]
    \item Largest team project (team size): [...]
\end{itemize}

\section{AI-Assisted Coding Policy}
% Specify to what extent the team plans to use generative AI, consistent with course policy:
% - AI only for work members could have done themselves given time/references
% - AI-generated/substantially-AI-designed code must be flagged in PR descriptions,
%   with source comments on non-trivial blocks
% - Team agrees on a common level of AI usage
% - Every PR requires review/approval from a member who understands the change and its source
[...]

\end{document}
